\documentstyle[%


righttag%


,11pt%


,amssymb%


,russian%               remove in Eng.


,izvru%                 Set izven% in Eng.


]{amsart}





% {\ge} {>}





\newcommand{\tts}[1]{}





%\hoffset=-15mm%                  For Eng.


\setcounter{page}{61}


\god{2005}


\nomer{10~(521)} %                        Remove in Eng.


%\nomer{Vol.\,49, No.\,10}%               Set in Eng.


%\str{\pageref{begin}--\pageref{end}}%  Set in Eng.


\udk{517.518}





\author{\it С.А.\,Степанянц}


% NB:   ^^ remove in Eng.


\title{Необходимые условия тауберова типа\\


для методов суммирования Чезаро}


\thanks{Работа выполнена при финансовой поддержке Российского фонда


фундаментальных исследований, грант \No\,03-01-00080.}





\date{}


%\pagestyle{myheadings}%         Set in Eng.


\pagestyle{plain} %              Remove in Eng.


\begin{document}


%\thispagestyle{plain}%          Set in Eng.


\maketitle


\label{begin}








{\bf 1.} В данной работе будут рассматриваться классические методы суммирования 


числовых рядов --- методы Чезаро ($C,\alpha$) с $\alpha \ge 0$


и метод Абеля ($A$) (определения этих методов см. ([1], сс.\,20,


125--128)), изучаться условия тауберова типа для этих методов.


 


Условимся, если не оговорено противное, везде далее считать, что 


${a_n} = \{a_n \}_{n=0}^{+\infty}$ --- последовательность


действительных чисел; $c_n=\{c_n \}_{n=0}^{+\infty}$ --- 


последовательность неотрицательных чисел; $\{n_r \}_{r=1}^{+\infty}$


--- последовательность натуральных чисел; 


$k,\ l,\ n$ --- целые неотрицательные числа; 


$\alpha,\ \beta$ --- неотрицательные числа. 


Запись $\sum a_n=A (P)$ означает, что ряд $\sum\limits_{n=0}^{+\infty}a_n$ 


суммируем к числу $A$ методом $P$ (ряды с неуказанными пределами 


суммирования будем считать


суммируемыми от 0 до $+\infty$), запись $a_n\in R$ означает, что 


последовательность $a_n$ удовлетворяет условию $R$.





Будем говорить, что метод $P$ включается методом $Q$ ($P\subset Q$), 


если из того, что $\sum a_n=A (P)$ следует $\sum a_n=A (Q)$.





Хорошо известны включения $(C,\alpha)\subset(C,\beta)\subset A$


при $\alpha < \beta$ ([1], сс.\,131, 140). 


Обратные включения в общем случае неверны,


однако они становятся справедливыми, если на ${a_n}$ наложить некоторые 


условия, называемые тауберовыми.





Будем рассматривать так называемые ``$O$-условия'' тауберова типа, т.\,е.


условия вида $a_n=O (c_n)$.





Скажем, что последовательность ${c_n}$ принадлежит классу $OT_Q (P)$


(${c_n} \in OT_Q (P)$), если из того,


что ряд $\sum a_n$ суммируем методом $P$, и того, что $a_n=O(c_n)$, следует, 


что $\sum a_n$ суммируем методом $Q$.


Метод $P$ при этом будем называть ``верхним'', $Q$ --- ``нижним''.


Далее в качестве ``верхнего'' метода рассматриваются методы Чезаро


$(C, \alpha)$ или Абеля $(A)$,


в качестве ``нижнего'' --- метод $(C,k)$ с $k<\alpha$.





Приведем некоторые классические результаты для случая, 


когда ``нижним'' методом является $(C, 0)$-сходимость.





Харди [2] доказал, что $\frac{1}{n} \in OT_{(C, 0)} (C, \alpha)$, 


Литлвуд [3] --- что $\frac{1}{n} \in OT_{(C, 0)} (A)$. 


 


В [4] рассмотрeно следующее условие $(E)$. 


Последовательность ${c_n} \in (E)$, если


для любого $\varepsilon>0$ существуют число $q>1$ и последовательность 


натуральных чисел $\{n_r\}_{r=1}^{\infty}$ такие, что


$\frac{n_{r+1}}{n_r} \ge q$


и 


$\sum\limits_{n_r<\nu<n_{r+1}} c_{\nu}<\varepsilon$ для всех $r$.





Пусть $P$ --- некоторый фиксированный метод $(C, \alpha)$ с $\alpha>0$ 


или метод $(A)$. Тогда, как показано в [4],


условие ${c_n} \in (E)$ является необходимым и достаточным для того, 


чтобы


$$


{c_n} \in OT_{(C, 0)} (P).


$$





Из этого результата следует, что условие Харди $a_n=O (\frac{1}{n})$ 


является наилучшим с точки зрения порядка, т.\,е. для любой 


последовательности ${w_n}$ такой, что 


$\lim\limits_{n \to+\infty} w_n=+\infty$, имеем


$$


{\frac{w_n}{n}} \notin OT_{(C, 0)} (P).


$$





Рассмотрим теперь аналогичные вопросы для случая, когда ``нижним'' методом 


является уже не метод $(C, 0)$, а метод $(C, k)$ с $k>0$.





В [5] установлено, что для любой последовательности $w_n$ такой, 


что $\lim\limits_{n \to+\infty} w_n=+\infty$,


существует $a_n$ такая, что $a_n=O(\frac{w_n}{n})$; ряд $\sum a_n$ 


суммируем методом $(A)$ и не суммируем никаким методом $(C, k)$, 


т.\,е. для всех $k$ имеем $\frac{w_n}{n} \notin OT_{(C, k)}(A)$.





В [6] были введены условия $(F, k)$, которые для краткости 


не будем здесь формулировать, и установлены следующие свойства этих условий:





1) $(F, 0)=(E)$;





2) eсли $k<l$ и $c_n\in(F,k)$, то $c_n\in(F,l)$;





3) eсли $k<l$, то существует $c_n$ такая, что $c_n\notin(F,k)$, но 


$c_n\in(F,l)$;





4) eсли $c_n\in(F,k)$, то $c_n\in OT_{(C, k)}(C,\alpha)$ для любого $\alpha>k$.


 


Из этих свойств следует, в частности, что условие $c_n \in (E)$ не является 


необходимым для того, чтобы $c_n \in OT_{(C, k)} (P)$, где $(P)$ --- метод


$(C,\alpha)$ с $\alpha>k$ или метод $(A)$. Ниже будут введены условия


$(ND, k)$ и получено, что требование $c_n \notin (ND, k)$


является необходимым для того, чтобы $c_n \in OT_{(C, k)} (P)$.





Доказательство этого проводится в теореме 1, которая формулируется в несколько 


иных обозначениях для придания ее утверждению более прозрачного смысла. 





Пусть $k$ --- натуральное число. Дадим определение условий $(ND, k)$. 





Последовательность $c_n$ удовлетворяет условию $(ND, k)$, 


если существуют числа $\varepsilon>0$ и $p$ ($0<p<1$) такие,


что для любого $q$ ($1<q<q^3 \le 2$) существуют следующие объекты, 


зависящие от $q$:





I) последовательность натуральных чисел $\{n_r^{(q)}\}_{r=1}^{+\infty}$ такая, 


что


$$


q \le \frac{n_{r+1}^{(q)}}{n_r^{(q)}} \le q^3;


$$





II) последовательность $\{\wt c{}_n^{(q)}\}$ такая, что $0\le\wt c{}_n^{(q)} 


\le c_n$;





III) $(k+2)$ возрастающих последовательностей 


натуральных чисел $\{l_{ir}^{(q)}\}_{r=1}^{+\infty}$ ($i=1,\dotsc,k+2$);





$(k+2)$ возрастающих последовательностей 


натуральных чисел $\{m_{ir}^{(q)}\}_{r=1}^{+\infty}$ ($i=1,\dotsc,k+2$), 


и для этих объектов для бесконечного числа номеров $r$ одновременно 


выполнены следующие соотношения:


\begin{gather}


n_r^{(q)} \le l_{1r}^{(q)}<m_{1r}^{(q)}<q^p m_{1r}^{(q)} \le 


l_{2r}^{(q)}<m_{2r}^{(q)}<q^p m_{2r}^{(q)} \le 


l_{3r}^{(q)}< \dotsb \le 


l_{(k+2)r}^{(q)}<m_{(k+2)r}^{(q)}<n_{r+1}^{(q)},\\


%


\sum_{\nu=l_{1r}^{(q)}+1}^{m_{1r}^{(q)}} \wt c{}_{\nu}^{(q)}=


\sum_{\nu=l_{ir}^{(q)}+1}^{m_{ir}^{(q)}} \wt c{}_{\nu}^{(q)} 


\quad (i=2, 3, \dotsc, k+2),\\


%


\sum_{i=1}^{k+2}\sum_{\nu=l_{ir}^{(q)}+1}^{m_{ir}^{(q)}} 


(n_{r+1}^{(q)}-\nu+1)^k \wt c{}_{\nu}^{(q)} \ge 


\varepsilon (n_r^{(q)})^k.


\end{gather}








\begin{thm}


Пусть $k$ --- фиксированное натуральное число, 


последовательность $c_n$ $(c_n {\ge} 0)$ удовлетворяет


условию $(ND, k)$. Тогда существует последовательность $a_n$, 


для которой выполнено


$1)$~$| a_n | \le c_n$ для всех $n;$ 


$2)$~ряд $\sum a_n$ суммируем к нулю любым методом $(C, \alpha)$ 


с $\alpha>k;$


$3)$~ряд $\sum a_n$ не суммируем методом $(C, k)$.


\end{thm}





Сделаем одно замечание об условиях $(ND, k)$. Эти условия были получены 


некоторым изменением (введением требования ``равномерности роста'') из


отрицания условий $(F, k)$ и внешне выглядят несколько громоздкими. Однако,


несмотря на это, они оказываются довольно удобными в применении. Так, простой


проверкой того, что ${\frac{w_n}{n}} \in (ND, k)$ для всех $k$, доказывается








\begin{thm}


Пусть последовательность ${w_n}$ такова, что


$\lim\limits_{n \to+\infty} w_n=+\infty$. 


Тогда $\frac{w_n}{n} \notin OT_{(C, k)}(C,\alpha)$ ни при каких


$k$ и $\alpha$   $(k<\alpha)$.


\end{thm}








Теорема 2 вместе с указанными выше результатами работ [2], [3], [5] полностью 


решает вопрос о наилучшем с точки зрения порядка $OT_{Q} (P)$-условии 


для случая $Q=(C, k)$; $P=(C, \alpha)$ с $\alpha>k$ или $P=A$, 


а именно $\frac{1}{n} \in OT_{(C, k)} (P)$; 


$\frac{w_n}{n} \notin OT_{(C, k)} (P)$, 


если $\lim\limits_{n \to+\infty} w_n=+\infty$. 


\medskip 





{\bf 2.} Установим некоторые вспомогательные утверждения, которые потребуются при


доказательстве теорем. Обозначения в приводимых ниже 


леммах максимально приближены к соответствующим обозначениям в теоремах.





В каждой из лемм фиксированы натуральные числа $k$, $l$, $m$, $l_i$, $m_i$, 


где $i=1,\dotsc,k+2$, причем выполнено


\begin{equation}


l \le l_1<m_1<l_2<m_2<\dotsb<l_{k+2}<m_{k+2} \le m.


\end{equation}


Действительные числа $\lambda$, $\beta_i$ ($i=1, \dotsc, k+2$), 


$\wt c_{\nu}$, $a_{\nu}$ ($\nu$ --- целое, $\nu \in (l; m]$)


заданы так, что 


\begin{alignat}{2}


\wt c_{\nu} &\ge 0 &\quad \text{для всех } \ \nu &\in (l; m ],\\


%


a_{\nu}&=0, &\quad \text{если } \ \nu &\in (l; m ] \setminus


\bigcup_{i=1}^{k+2} (l_i; m_i],\\


%


a_{\nu}&=\lambda \cdot \beta_i \cdot


\wt c_{\nu}, &\quad \text{если } \ \nu &\in (l_i; m_i].


\end{alignat}


Будут использоваться также следующие обозначения: $[t]$ --- целая часть 


числа $t$;


$$


Y_{\nu}(n, k)=\prod\limits_{s=1}^{k} (n-\nu+s).


$$








\begin{lem}


Пусть $k$, $l$, $m$, $l_i$, $m_i$ $(i=1, \dotsc, k+2)$ --- натуральные числа, 


удовлетворяющие неравенствам $(4)$. Пусть $\lambda$, $\beta_i$


$(i=1, \dotsc, k+2)$, $\wt c_{\nu}$, $a_{\nu}$ 


$(\nu$ --- целое, $\nu\in(l; m])$ --- действительные числа,


удовлетворяющие соотношениям $(5)$--$(7)$ 


и такие, что для всех целых $t \in [0; k]$ выполнено


\begin{equation}


\sum_{i=1}^{k+2} \bigg(\sum_{\nu=l_i+1}^{m_i} \nu^t \wt c_{\nu}\bigg)\beta_i=0.


\end{equation}


Пусть $n$ --- натуральное число такое, что $n \ge m$.


Тогда справедливы равенства


\begin{gather}


\sum_{\nu=l+1}^m Y_{\nu}(n, k+1) a_{\nu}=(-1)^{k+1} \lambda


\bigg(\sum_{i=1}^{k+2} \bigg(\sum_{\nu=l_i+1}^{m_i} \nu^{k+1}


\wt c_{\nu}\bigg) \beta_i\bigg)=


% 


\lambda \sum_{i=1}^{k+2} \beta_i\bigg(\sum_{\nu=l_i+1}^{m_i}


(m-\nu+1)^{k+1} \wt c_{\nu}\bigg),\\


%


\sum_{\nu=l+1}^m Y_{\nu}(n, k) a_{\nu}=0.


\end{gather}


\end{lem}








\begin{pf}


Выражение $Y_{\nu} (n, k+1)$ является многочленом степени $k+1$ от 


переменной~$\nu$. Обозначим числовой коэффициент при $\nu^t$ в стандартной 


записи этого многочлена через~$\chi_t$, т.\,е. 


$Y_{\nu} (n, k+1)=\sum\limits_{t=0}^{k+1} \chi_t \nu^t$.


Тогда, т.\,к. $\chi_{k+1}=(-1)^{k+1}$, то, используя (5)--(8), получаем


\begin{multline*}


\sum\limits_{\nu=l+1}^m Y_{\nu}(n, k+1)a_{\nu}=\lambda


\sum\limits_{i=1}^{k+2}\beta_i\sum\limits_{\nu=l_i+1}^{m_i}


\bigg(\sum\limits_{t=0}^{k+1} \chi_t\nu^t\bigg) \wt c_{\nu}= \\


%


=\lambda\sum\limits_{t=0}^{k+1}\chi_t\bigg(\sum\limits_{i=1}^{k+2}


\bigg(\sum\limits_{\nu=l_i+1}^{m_i} \nu^t 


\wt c_{\nu}\bigg)\beta_i\bigg)=(-1)^{k+1}\lambda


\bigg(\sum\limits_{i=1}^{k+2}\bigg(\sum\limits_{\nu=l_i+1}^{m_i} 


\nu^{k+1} \wt c_{\nu}\bigg) \beta_i\bigg). 


\end{multline*}





Аналогично, раскрывая по степеням $\nu$ выражения $(m-\nu+1)^k$ и 


$Y_{\nu} (n, k)$ и применяя (8), получаем соответственно второе из 


равенств формулы (9) и формулу (10). 


\end{pf}





При доказательстве леммы 2 будут использоваться следующие неравенства, 


справедливые для всех $q \in (1;\, 1{,}5)$,   %%% <-- use {.} in eng.


$p \in (0; 1)$


\begin{align}


\frac{q^p-1}{q-1} &\ge \frac{p}{2p+2},\\


q^8-1&<50(q-1).


\end{align}


 


Неравенство (11) вытекает из возрастания на $[0;1]$ функции


$\phi(y)=(2y+2)(q^y -1)-y(q -1)$, 


которое, в свою очередь, следует из того, что $\phi'(0)>0$ и 


$\phi''(y)>0$ при $y\in [0;1]$. Неравенство (12) получается 


разложением $q^8-1$ на множители и элементарными числовыми оценками.








\begin{lem}


Пусть $q$ $(1<q<q^3<2)$, $p$ $(0<p<1)$ ---


фиксированные числа\/$;$ $k$, $l$, $m$, $l_i$, $m_i$ $(i=1, \dotsc, k+2)$ --- 


натуральные числа, удовлетворяющие неравенствам $(4)$ и такие, что


\begin{equation}


q^p m_i \le l_{i+1}, \quad q \le \frac{m}{l} \le q^3,


\quad 1+\frac{k+1}{m}<q^3.


\end{equation}





Пусть $l_{k+3}=[q^p m]+1$, $V$ --- фиксированное положительное число,


$\lambda$ $(\lambda{>}0)$, $\beta_i$ $(i{=}1,..., k{+}2)$,


%%%%%%%%% remove extra { }, put \dotsc       ^^^^


$\wt c_{\nu}$, $a_{\nu}$ $(\nu$ --- целое,


$\nu \in (l; m])$ --- действительные числа, удовлетворяющие соотношениям


$(5)$--$(7)$ и равенству


$\sum\limits_{\nu=l_1+1}^{m_1}\wt c_{\nu}=\sum\limits_{\nu=l_i+1}^{m_i}


\wt c_{\nu}$ $(i=2, \dotsc, k+1)$. 





Пусть справедливы неравенства


\begin{equation}


\bigg|\sum_{\nu=l+1}^{m_i} Y_{\nu}(l_{i+1}; k) a_{\nu} \bigg|<


\frac{V}{2^{k+2}} \bigg(\frac{p}{100p+100}\bigg)^{k(k+2)} m_i^k


\end{equation}


для $i=1, \dotsc, k+2$ и


\begin{equation}


\lambda \sum_{\nu=l_1+1}^{m_1} \wt c_{\nu}>\frac{V}{(q^8-1)^k}.


\end{equation}





Тогда для всех целых $i \in [1; k+2]$ имеем


\begin{equation}


|\beta_i|<\frac{1}{2^{k+3-i}}\bigg(\frac{p}{100p+100}\bigg)^{k(k+2-i)}. 


\end{equation}


\end{lem}








{\bf Доказательство} проведем индукцией по $i$. 


Пусть $i=1$. Тогда, применяя неравенства (11)--(15), получаем


\begin{multline*}


\frac{V}{2^{k+2}} \bigg(\frac{p}{100p+100}\bigg)^{k(k+2)} m_i^k>


\bigg|\sum_{\nu=l+1}^{m_1} Y_{\nu} (l_2; k) a_{\nu} \bigg|=


\lambda |\beta_1| \sum_{\nu=l_1+1}^{m_1} Y_{\nu} (l_2; k) 


\wt c_{\nu} \ge \\


%


\ge \lambda |\beta_1| (l_2-m_1)^k 


\sum_{\nu=l_1+1}^{m_1} \wt c_{\nu} \ge 


\frac{V}{50^k (q-1)^k} |\beta_1| m_1^k \bigg(\frac{l_2}{m_1}-1\bigg)^k \ge \\


%


\ge \frac{V (q^p-1)^k}{50^k (q-1)^k} | \beta_1 | m_1^k \ge 


V \bigg(\frac{p}{100p+100}\bigg)^k |\beta_1| m_1^k. 


\end{multline*}


Отсюда следует, что


$|\beta_1|<\frac{1}{2^{k+2}} (\frac{p}{100p+100})^{k(k+1)}$. 


Для $i=1$ неравенство (16) доказано.





Допустим, что оно верно для $i=1,\dotsc,w$, где $w$ --- целое,


$w \in [1, k+1]$. Покажем, что оно верно для $i=w+1$. 


С использованием (13), (14) имеем


\begin{multline*}


\frac{V}{2^{k+2}} \bigg(\frac{p}{100p+100}\bigg)^{k(k+2)} m_{w+1}^k>


\bigg| \sum_{\nu=l+1}^{m_{w+1}} Y_{\nu}(l_{w+2}; k) a_{\nu} \bigg|=\\


%


=\bigg|\lambda\sum_{i=1}^w\beta_i\sum_{\nu=l_i+1}^{m_{i}} Y_{\nu}(l_{w+2}; k) 


\wt c_{\nu}+\lambda \beta_{w+1} \sum_{\nu=l_{w+1}+1}^{m_{w+1}} 


Y_{\nu}(l_{w+2}; k) \wt c_{\nu} \bigg| \ge \\


%


\ge \lambda |\beta_{w+1}|\sum_{\nu=l_{w+1}+1}^{m_{w+1}} Y_{\nu}(l_{w+2};k)


\wt c_{\nu}-\lambda \sum_{i=1}^{w} | \beta_i | 


\sum_{\nu=l_i+1}^{m_{i}} Y_{\nu}(l_{w+2}; k) \wt c_{\nu} \ge \\


%


\ge \lambda | \beta_{w+1} | (l_{w+2}-m_{w+1})^k \sum_{\nu=l_{w+1}+1}^


{m_{w+1}} \wt c_{\nu}-\lambda (l_{w+2}-l_1+k)^k \sum_{i=1}^w | 


\beta_i | \sum_{\nu=l_{i}+1}^{m_{i}} \wt c_{\nu}= \\


%


= \lambda \bigg(\sum_{\nu=l_1+1}^{m_1} \wt c_{\nu}\bigg)\bigg(| \beta_{w+1} | 


m_{w+1}^k \bigg(\frac{l_{w+2}}{m_{w+1}} -1\bigg)^k-


\bigg(\sum_{i=1}^w | \beta_i |\bigg) m_{w+1}^k


\bigg(\frac{l_1}{m_{w+1}}\bigg)^k


\bigg(\frac{l_{w+2}+k}{l_1}-1\bigg)^k\bigg) \ge \\


%


\ge m_{w+1}^k \lambda \bigg(\sum_{\nu=l_1+1}^{m_1} \wt c_{\nu}\bigg)


\bigg(| \beta_{w+1} | 


(q^p-1)^k-\bigg(\sum_{i=1}^w | \beta_i |\bigg) (q^8-1)^k\bigg).


\end{multline*}


Последний переход совершен с помощью неравенства


$\frac{l_{w+2}+k}{l_1}=\frac{(q^p m+1+k) q^p m}{q^p m l_1} 


<(1+\frac{1+k}{m}) q^p \frac{m}{l}<q^3 q q^3<q^8$.


Применяя теперь неравенство (15), получаем


$$


\frac{V}{2^{k+2}} \bigg(\frac{p}{100p+100}\bigg)^{k(k+2)} \ge 


V \bigg(| \beta_{w+1} |


\bigg(\frac{q^p -1}{q^8-1}\bigg)^k-\sum_{i=1}^{w} | \beta_i |\bigg).


$$


Отсюда, используя формулы (11), (12) и индукционное предположение, заключаем


\begin{multline*}


| \beta_{w+1} |<\bigg(\frac{q^8 -1}{q^p-1}\bigg)^k 


\bigg(\sum_{i=1}^{w} | \beta_i |+\frac{1}{2^{k+2}}


\bigg(\frac{p}{100p+100}\bigg)^{k(k+2)}\bigg)< \\


%


< \bigg(\frac{100p+100}{p}\bigg)^k \bigg(\sum_{i=1}^{w} \frac{1}{2^{k+3-i}} 


\bigg(\frac{p}{100p+100}\bigg)^{k(k+2 -i)}+\frac{1}{2^{k+2}}


\bigg(\frac{p}{100p+100}\bigg)^{k(k+2)}\bigg)< \\


%


< \bigg(\frac{100p+100}{p}\bigg)^k \bigg(\frac{p}{100p+100}\bigg)^{k(k+2-w)} 


\bigg(\sum_{i=1}^{w} \frac{1}{2^{k+3-i}}+\frac{1}{2^{k+2}}\bigg)=


%


\bigg(\frac{p}{100p+100}\bigg)^{k(k+1-w)} \frac{1}{2^{k+2-w}}.\quad\square


\end{multline*}








{\bf 3. Доказательство теоремы 1.} Так как $c_n\in(ND,k)$, то существуют 


$\varepsilon>0$ и $p$ ($0<p<1$) такие, что для любого $q$ ($1<q<q^3 \le 2$) 


найдутся указанные в условии $(ND, k)$ объекты, зависящие, 


вообще говоря, от $q$.





В качестве $q$ будем рассматривать величины


$q_s=\sqrt[\uproot{2}3]{1+(\frac{1}{2})^{s+2}}$ при натуральных $s$.





Для каждого из этих $q_s$ существует следующий набор ``своих'' объектов


(в качестве верхнего индекса у них


правильно было бы писать $(q_s)$, но для краткости будем писать $(s)$):


\begin{itemize}


\item[I)] последовательность натуральных чисел $\{n_r^{(s)}\}_{r=1}^{+\infty}$


такая, что 


$$


q_s \le \frac{n_{r+1}^{(s)}}{n_r^{(s)}} \le q_s^3;


$$


\item[II)] последовательность $\{\wt c{}_n^{(s)} \}_{n=0}^{+\infty}$ такая, что


$0 \le \wt c{}_n^{(s)} \le c_n$;


\item[III)] $(2k+4)$ возрастающих последовательностей натуральных чисел


$\{l_{ir}^{(s)} \}_{r=1}^{+\infty}$, 


$\{m_{ir}^{(s)}\}_{r=1}^{+\infty}$\linebreak ($i=1,\dotsc, k+2$),


\end{itemize}


и для этих объектов для бесконечного числа номеров $r$ 


одновременно выполнены соотношения (1)--(3).





Построим теперь по индукции последовательности 


$\{l^{(s)} \}_{s=1}^{+\infty}$, $\{m^{(s)} \}_{s=1}^{+\infty}$,


$\{l_i^{(s)} \}_{s=1}^{+\infty}$, $\{m_i^{(s)} \}_{s=1}^{+\infty}$\linebreak


($i=1, \dotsc, k+2$).


Для этого обозначим через ${\rho_1}$ первый из номеров $r$ 


последовательности  $\{n_r^{(1)}\}$, для которых выполнены (1)--(3) и 


$n_{\rho_1}^{(1)}>64 (k+1)$.





Положим


$$


l^{(1)}=n_{\rho_1}^{(1)}, \ \ m^{(1)}=n_{\rho_1+1}^{(1)},


\ \ l_i^{(1)}=l_{i \rho_1}^{(1)},\ \ m_i^{(1)}=m_{i \rho_1}^{(1)}


\quad (i=1, \dotsc, k+2).


$$


Тем самым первые члены последовательностей построены. Пусть теперь построены 


первые\linebreak $(s-1)$ членов. Построим члены с номерами $s$.





Обозначим через $\rho_s$ первый из номеров последовательности


$\{n_r^{(s)} \}$, для которых выполнены соотношения (1)--(3)


и справедливы следующие неравенства: 


\begin{gather*}


n_{\rho_s}^{(s)}>2m^{(s-1)}, \\


\sum_{j=1}^{s-1}(l^{(j)})^{k+1}\Big/ (n_{\rho_s}^{(s)})^{k+1}<\frac{1}{2^{s+1}}. 


\end{gather*}


Положим $l^{(s)}=n_{\rho_s}^{(s)}$, $m^{(s)}=n_{\rho_s+1}^{(s)}$,


$l_i^{(s)}=l_{i \rho_s}^{(s)}$, $m_i^{(s)}=m_{i \rho_s}^{(s)}$.





Тем самым последовательности $\{l^{(s)} \}$, $\{m^{(s)} \}$, 


$\{l_i^{(s)} \}$, $\{m_i^{(s)} \}$, где $i=1,\dotsc, k+2$, $s \in \Bbb N$,


построены, причем для всех $s$ выполнены соотношения


\begin{gather}


1<q_s \le \frac{m^{(s)}}{l^{(s)}} \le q_s^3=1+\frac{1}{2^{s+2}},\\


%


8(k+1)<l^{(s)}\le l_1^{(s)}<m_1^{(s)}<


q_s^p m_1^{(s)}\le l_2^{(s)}<m_2^{(s)}<q_s^p m_2^{(s)} \le 


%


l_3^{(s)}< \dotsb 


\le l_{k+2}^{(s)}<m_{k+2}^{(s)} \le m^{(s)},\\


%


\sum_{\nu=l_1^{(s)}+1}^{m_1^{(s)}} \wt c{}_{\nu}^{(s)}=


\sum_{\nu=l_i^{(s)}+1}^{m_i^{(s)}} \wt c{}_{\nu}^{(s)}\quad (i=2,\dotsc,k+2),\\


%


\sum_{i=1}^{k+2} \sum_{\nu=l_i^{(s)}+1}^{m_i^{(s)}} (m^{(s)}-\nu+1)^k 


\wt c{}_{\nu}^{(s)} \ge \varepsilon (l^{(s)})^k,\\


%


l^{(s)}>2 m^{(s-1)},\\


%


\frac{\sum\limits_{j=1}^{s-1} (l^{(j)})^{k+1}}{(l^{(s)})^{k+1}}


<\frac{1}{2^{s+1}}.


\end{gather}





Из соотношений (18) и (20) получим, что последовательности 


$\{l^{(s)}\}$ и $\{m^{(s)}\}$ монотонно возрастают, стремясь к $+\infty$,


и полуинтервалы $(l^{(s)}; m^{(s)}]$ при различных $s$ попарно не пересекаются.





Далее будем строить последовательность $\{a_n \}_{n=0}^{+\infty}$,


удовлетворяющую условиям теоремы. 





Положим $a_n=0$, если


$n \notin \mathop{\cup}\limits_{s=1}^{+\infty}(l^{(s)};m^{(s)}]$.


Будем считать теперь, что $s$ --- произвольное, но фиксированное натуральное 


число, и построим $a_n$ для $n \in (l^{(s)}; m^{(s)}]$.


Для этого понадобятся числа $\lambda^{(s)}$ и $\beta_i^{(s)}$, где 


$i$ --- целое, $i \in [1, k+2]$, которые сейчас определим.


Положим 


\begin{equation}


\lambda^{(s)}=\frac{\varepsilon (l^{(s)})^k}


{\sum\limits_{i=1}^{k+2} \sum\limits_{\nu=l_i^{(s)}+1}^{m_i^{(s)}} 


(m^{(s)}-\nu+1)^k \wt c{}_{\nu}^{(s)}}.


\end{equation}


Из (20) следует, что $0<\lambda^{(s)} \le 1$. 





Рассмотрим теперь следующую систему линейных уравнений относительно 


неизвестных $\beta_i^{(s)}$, где $i$ --- целое, $i \in [1, k+2]$:


\begin{equation}


\begin{split}


\bigg(\sum_{\nu=l_1^{(s)}+1}^{m_1^{(s)}} \wt c{}_{\nu}^{(s)}\bigg)


\beta_1^{(s)}+\dotsb+\bigg(\sum_{\nu=l_{k+2}^{(s)}+1}^{m_{k+2}^{(s)}}


\wt c{}_{\nu}^{(s)}\bigg) \beta_{k+2}^{(s)}&=0,\\


%


\bigg(\sum_{\nu=l_1^{(s)}+1}^{m_1^{(s)}} \nu \wt c{}_{\nu}^{(s)}\bigg)


\beta_1^{(s)}+\dotsb+\bigg(\sum_{\nu=l_{k+2}^{(s)}+1}^{m_{k+2}^{(s)}}


\nu \wt c{}_{\nu}^{(s)}\bigg) \beta_{k+2}^{(s)}&=0, \\


%


\dotsc\dotsc\dotsc& \phantom{=}\\


%


\bigg(\sum_{\nu=l_1^{(s)}+1}^{m_1^{(s)}} \nu^k \wt c{}_{\nu}^{(s)}\bigg)


\beta_1^{(s)}+\dotsb+\bigg(\sum_{\nu=l_{k+2}^{(s)}+1}^{m_{k+2}^{(s)}}


\nu^k \wt c{}_{\nu}^{(s)}\bigg) \beta_{k+2}^{(s)}&=0.


\end{split}


\end{equation}





В системе (24) количество уравнений $(k+1)$ меньше числа неизвестных.


Следовательно, эта система обладает решениями, отличными от нулевого. 


Тогда, очевидно, существует решение 


$(\beta_1^{(s)},\beta_2^{(s)}, \dotsc, \beta_{k+2}^{(s)})$ 


системы (24) такое, что $\max\limits_{1 \le i \le k+2}|\beta_i^{(s)}|=1$.


Именно такой конкретный набор чисел $\beta_i^{(s)}$ и будем в дальнейшем 


рассматривать.





Положим 


\begin{alignat*}{2}


a_n&=\lambda^{(s)} \beta_i^{(s)} c_n^{(s)} &\quad \text{при } 


\ n &\in (l_i^{(s)}; m_i^{(s)}],\ \text{ где }\ i=1, \dotsc, k+2;\\


%


a_n&=0 &\quad \text{при } \ n &\in (l^{(s)}; m^{(s)}]


\setminus \mathop{\cup}\limits_{i=1}^{k+2} (l_i^{(s)}; m_i^{(s)}].


\end{alignat*} 





Таким образом, последовательность $\{a_n \}_{n=0}^{+\infty}$ определена, 


причем $|a_n|\le c_n$ для всех $n$. 





Покажем, что $\sum a_n=0 \ (C, k+1)$.


Возьмем любое $\delta>0$. Для этого $\delta$ существует натуральное


$s_0$ ($s_0>2$) такое, что для всех $s \ge s_0$ выполнено неравенство 


\begin{equation}


\frac{(k+1)!}{2^s} \varepsilon<\delta.


\end{equation}





Пусть $n>l^{(s_0)}$. Тогда существует $s \ge s_0$ такое, что


$l^{(s)}<n \le l^{(s+1)}$.


Обозначим $T^k (n)=\sum\limits_{\nu=0}^n Y_{\nu} (n, k) a_{\nu}$, 


$T^{k+1} (n)=\sum\limits_{\nu=0}^n Y_{\nu} (n, k+1) a_{\nu}$


и оценим модули этих величин.





Применяя лемму 1 к каждому из промежутков


$(l^{(j)};m^{(j)}]$, где $j=1,\dotsc,s-1$, т.\,е. полагая


в лемме 1 \ $l=l^{(j)}$, $m=m^{(j)}$, $l_i=l_i^{(j)}$, $m_i=m_i^{(j)}$, 


$\lambda=\lambda^{(j)}$, $c_{\nu}=c_{\nu}^{(j)}$, получаем


\begin{gather}


|T^{k+1} (n)|=\bigg| \sum_{j=1}^{s-1} 


\sum_{\nu=l^{(j)}+1}^{m^{(j)}} Y_{\nu} (n, k+1) a_{\nu}+


\sum_{\nu=l^{(s)}+1}^n Y_{\nu} (n, k+1) a_{\nu} \bigg|=\notag \\


%


=\bigg| \sum_{j=1}^{s-1} \lambda^{(j)} \sum_{i=1}^{k+2} \beta_i 


\bigg(\sum_{\nu=l_i^{(j)}+1}^{m_i^{(j)}} (m^{(j)}-\nu+1)^{k+1} 


\wt c{}_{\nu}^{(j)}\bigg)+


\sum_{\nu=l^{(s)}+1}^n Y_{\nu} (n, k+1) a_{\nu} \bigg|;\notag\\


%


|T^k (n)|=\bigg|\sum_{\nu=l^{(s)}+1}^n Y_{\nu} (n, k) a_{\nu} \bigg|.


\end{gather}





Рассмотрим отдельно два случая расположения $n$ в промежутке 


$(l^{(s)}; l^{(s+1)}]$:





A) $m^{(s)}<n \le l^{(s+1)}$,





Б) $l^{(s)}<n \le m^{(s)}$.





В случае A), еще раз применяя лемму 1 к промежутку 


$(l^{(s)}; m^{(s)}]$, выводим 


\begin{gather}


|T^{k+1} (n)|=\bigg| \sum_{j=1}^{s} \lambda^{(j)} 


\sum_{i=1}^{k+2} \beta_i 


\sum_{\nu=l_i^{(j)}+1}^{m_i^{(j)}} (m^{(j)}-\nu+1)^{k+1} 


\wt c{}_{\nu}^{(j)} \bigg| \le 


%


\sum_{j=1}^{s} \lambda^{(j)} \sum_{i=1}^{k+2} 


\sum_{\nu=l_i^{(j)}+1}^{m_i^{(j)}} (m^{(j)}-\nu+1)^{k+1} 


\wt c{}_{\nu}^{(j)},\\


%


|T^k (n) |=0.


\end{gather}


В случае Б), используя оценки


$Y_{\nu} (n, k)<k! (m^{(s)}-\nu+1)^{k}$, 


$Y_{\nu} (n, k+1)<(k+1)! (m^{(s)}-\nu+1)^{k+1}$, получаем 


\begin{align}


|T^{k+1} (n)| &\le (k+1)! \sum_{j=1}^{s} \lambda^{(j)} 


\sum_{i=1}^{k+2} \sum_{\nu=l_i^{(j)}+1}^{m_i^{(j)}} (m^{(j)}-\nu+1)^{k+1} 


\wt c{}_{\nu}^{(j)},\\


%


| T^{k} (n) | &\le k! \lambda^{(s)} \sum_{i=1}^{k+2} 


\sum_{\nu=l_i^{(s)}+1}^{m_i^{(s)}} (m^{(s)}-\nu+1)^{k} 


\wt c{}_{\nu}^{(s)}.


\end{align}


Сравнивая полученные оценки с формулами (27), (28), заключаем, что для любого 


$n \in (l^{(s)}; l^{(s+1)}]$ верны неравенства (29), (30).





Из формулы (29) с использованием (17), (22), (23), (25) получаем


\begin{multline}


\frac{| T^{k+1} (n) |}{(n+1) \dotsb (n+k+1)} \le 


\frac{(k+1)!}{(l^{(s)})^{k+1}} \sum_{j=1}^{s} 


\lambda^{(j)} l^{(j)} \bigg(\frac{m^{(j)}}{l^{(j)}}-1\bigg)\sum_{i=1}^{k+2} 


\sum_{\nu=l_i^{(j)}+1}^{m_i^{(j)}} (m^{(j)}-\nu+1)^{k}\wt c{}_{\nu}^{(j)} \le \\


%


\le \frac{(k+1)!}{(l^{(s)})^{k+1}}\sum_{j=1}^{s} \frac{\varepsilon 


(l^{(j)})^{k+1}} {2^{j+2}}<\frac{(k+1)! \varepsilon}{(l^{(s)})^{k+1}}


\bigg(\sum_{j=1}^{s-1}(l^{(j)})^{k+1}+


\frac{(l^{(s)})^{k+1}}{2^{s+2}} \bigg)< 


%


\frac{(k+1)!}{2^s} \varepsilon<\delta.


\end{multline}


Формула (31) означает, что 


$$


\lim_{n\to+\infty}\frac{1}{{n+k+1 \choose k+1}}\sum_{\nu=0}^n


{n-\nu+k+1 \choose k+1} a_{\nu}=0.


$$


Следовательно ([1], с.\,125--126), $\sum a_n=0 \ (C, k+1)$.





Из формул (23) и (30) сразу вытекает, что


$|T^{k} (n)|\le k! \varepsilon n^k$ и, следовательно, 


$\sum\limits_{\nu=0}^n {n-\nu+k \choose k} a_{\nu}=O(n^k)$, 


а это по определению ([1], с.\,128) означает, что 


ряд $\sum a_n$ ограничен $(C, k)$.





Таким образом, ряд $\sum a_n$ ограничен $(C, k)$ и суммируем к нулю 


методом $(C, k+1)$. Следовательно, по теореме о выпуклости ([1], с.\,163--164,


теорема 70) $\sum a_n=0 \ (C, \alpha)$ для всех $\alpha>k$.





Для доказательства теоремы осталось установить, что 


ряд $\sum a_n$ не суммируем методом $(C, k)$. Предположим противное, что


ряд $\sum a_n$ суммируем $(C, k)$.


Поскольку $\sum a_n=0 \ (C, k+1)$, то и суммироваться $(C, k)$ 


ряд $\sum a_n$ мог бы только к числу~0. Тогда


$$


\lim_{n \to+\infty}\frac{1}{{n+k \choose k}}


\sum_{\nu=0}^n {n-\nu+k \choose k} a_{\nu}=0.


$$


Это соотношение эквивалентно тому, что


$\sum\limits_{\nu=0}^n Y_{\nu} (n, k) a_{\nu}=o(n^k)$.


Тогда для любого $B$ существует~$s$ такое, что для любого


$n>l^{(s)}$ имеем 


$\frac{1}{n^k}\Big| \sum\limits_{\nu=0}^n Y_{\nu} (n, k) a_{\nu} \Big|<B$.





В качестве $B$ будем рассматривать величину 


$\frac{\varepsilon}{(k+2) 2^{2k+2}}(\frac{p}{100p+100} )^{k(k+2)}$. 


Тогда из (26) для\linebreak $i=1,\dotsc,k+1$ получаем


\begin{multline*}


\bigg|\sum_{\nu=0}^{l_{i+1}^{(s)}} Y_{\nu} (l_{i+1}^{(s)}, k) a_{\nu} \bigg|=


|T^k (l_{i+1}^{(s)})|=\bigg|\sum_{\nu=l^{(s)}+1}^{l_{i+1}^{(s)}}


Y_{\nu} (l_{i+1}^{(s)}, k) a_{\nu} \bigg|=\\


%


=\bigg|\sum_{\nu=l^{(s)}+1}^{m_{i}^{(s)}}


Y_{\nu} (l_{i+1}^{(s)}, k) a_{\nu} \bigg|<


B (l_{i+1}^{(s)})^k<2^k B (m_{i}^{(s)})^k.


\end{multline*}


Положив $l_{k+3}^{(s)}=[q^p m^{(s)}]+1$, по лемме 1 получим 


$$


\bigg|\sum_{\nu=l^{(s)}+1}^{m_{k+2}^{(s)}}


Y_{\nu} (l_{k+3}^{(s)}, k) a_{\nu} \bigg|=


\bigg|\sum_{\nu=l^{(s)}+1}^{m^{(s)}}


Y_{\nu} (l_{k+3}^{(\nu)}, k) a_{\nu} \bigg|=0.


$$ 


Тем самым для всех $i=1,\dotsc, k+2$ имеем


\begin{equation}


\bigg|\sum_{\nu=l^{(s)}+1}^{m_{i}^{(s)}}


Y_{\nu} (l_{i+1}^{(s)}, k) a_{\nu} \bigg|<


\frac{\varepsilon}{(k+2) 2^{(k+2)}} 


\bigg(\frac{p}{100p+100} \bigg)^{k(k+2)} (m_{i}^{(s)})^k.


\end{equation}


Далее 


\begin{gather*}


\varepsilon (l^{(s)})^k=


\lambda^{(s)} \sum_{i=1}^{k+2} \sum_{\nu=l_i^{(s)}+1}^{m_i^{(s)}} 


(m^{(s)}-\nu+1)^k \wt c{}_{\nu}^{(s)} \le 


\lambda^{(s)} \sum_{i=1}^{k+2} (m^{(s)}-l^{(s)})^k 


\sum_{\nu=l_i^{(s)}+1}^{m_i^{(s)}} \wt c{}_{\nu}^{(s)}= \\


%


= (k+2) \lambda^{(s)} (l^{(s)})^k \bigg(\frac{m^{(s)}}{l^{(s)}}-1\bigg)^k 


\sum_{\nu=l_1^{(s)}+1}^{m_1^{(s)}} \wt c{}_{\nu}^{(s)}<


(k+2) (l^{(s)})^k (q^8-1)^k 


\lambda^{(s)} \sum_{\nu=l_1^{(s)}+1}^{m_1^{(s)}} \wt c{}_{\nu}^{(s)}.


\end{gather*}


Следовательно,


\begin{equation}


\lambda^{(s)} \sum_{\nu=l_1^{(s)}+1}^{m_1^{(s)}} \wt c{}_{\nu}^{(s)}>


\frac{\varepsilon}{(k+2)(q^8-1)^k}.


\end{equation}





Сравнивая формулы (32), (33) с формулами (14), (15), замечаем, что выполнены 


условия леммы 2 на промежутке $ (l^{(s)}; m^{(s)}]$


с


$V=\frac{\varepsilon}{k+2}$


(соотношение $1+\frac{k+1}{m^{(s)}}<q^3$ выполнено, т.\,к. из (21) следует


$\frac{k+1}{m^{(s)}}<\frac{k+1}{l^{(1)}} \frac{l^{(1)}}{l^{(2)}} \dotsb 


\frac{l^{(s-1)}}{l^{(s)}}<\frac{1}{8}


\frac{1}{2^{s-1}}=\frac{1}{2^{s+2}}=q_s^3-1$). 





Применяя лемму 2, получаем для всех целых $i \in [1, k+2]$ неравенство


$$


|\beta_i^{(s)}|<\frac{1}{2^{k+3-i}}\bigg(\frac{p}{100p+100}\bigg)^{k(k+2 -i)}


\le \frac{1}{2},


$$ 


которое противоречит тому, что


$\max\limits_{1 \le i \le k+2} | \beta_i^{(s)} |=1$.


Полученное противоречие доказывает теорему~1.$\quad\square$ 


\medskip





{\bf 4. Доказательство теоремы 2.}


Пусть $c_n=\frac{w_n}{n}$ для всех $n$ и пусть $k$ --- любое фиксированное 


натуральное число. Согласно теореме 1 достаточно установить, что 


$c_n \in (ND, k)$.





Возьмем $\varepsilon=1$, $p=\frac{1}{2k+4}$ и построим для любого $q$ 


такого, что $1<q<q^3 \le 2$, требуемые в условии $(ND, k)$ объекты.


Начнем с построения последовательностей


$\{n_r^{(q)}\}_{r=1}^{+\infty}$, $\{l_{ir}^{(q)} \}_{r=1}^{+\infty}$, 


$\{m_{ir}^{(q)} \}_{r=1}^{+\infty}$ ($i=1, \dotsc, k+2$).


Положим $n_1^{(q)}=\big[\frac{1}{q^{2p}-1}\big]+1$.


Считая теперь, что при некотором натуральном $r$ построена величина 


$n_r^{(q)}$, определим значения


$l_{ir}^{(q)}$, $m_{ir}^{(q)}$ ($i=1,\dotsc,k+2$)


и


$n_{r+1}^{(q)}$ следующим образом:


$l_{1r}^{(q)}=n_r^{(q)}$, $m_{ir}^{(q)}=[q^p l_{ir}^{(q)}]+1$ 


($i=1,\dotsc, k+2$), $l_{(i+1)r}^{(q)}=[q^p m_{ir}^{(q)}]+1$ 


($i=1, \dotsc, k+1$), $n_{r+1}^{(q)}=[q^p m_{(k+2)r}^{(q)}]+1$.





Тем самым требуемые последовательности построены и для них выполнены 


соотношения


\begin{equation}


\frac{n_{r+1}^{(q)}}{n_{r}^{(q)}}=


\frac{n_{r+1}^{(q)}}{m_{(k+2)r}^{(q)}} 


\prod_{i=1}^{k+2} \frac{m_{ir}^{(q)}}{l_{ir}^{(q)}} 


\prod_{i=1}^{k+1} \frac{l_{(i+1)r}^{(q)}}{m_{ir}^{(q)}}>


q^p (q^p)^{k+2} (q^p)^{k+1}=q,


\end{equation}


$\frac{m_{ir}^{(q)}}{l_{ir}^{(q)}} \le q^p+\frac{1}{l_{ir}^{(q)}} \le 


q^p+\frac{1}{n_{1}^{(q)}}<q^p+(q^{2p}-1)<


q^p+q^p (q^{2p} -1)=q^{3p}$ для всех целых $i\in [1, k+2]$.





Аналогично 


$\frac{l_{(i+1)r}^{(q)}}{m_{ir}^{(q)}}<q^{3p}$ ($i=1, \dotsc, k+1$)


и


$\frac{n_{r+1}^{(q)}}{m_{(k+2)r}^{(q)}}<q^{3p}$.





Далее, применяя разложение $\frac{n_{r+1}^{(q)}}{n_{r}^{(q)}}$ 


в произведение (см. (34)), получаем


$\frac{n_{r+1}^{(q)}}{n_{r}^{(q)}}<q^3$. 


Кроме того, из построения следует справедливость формулы (1) 


из определения $(ND, k)$.





Построим последовательность $\{\wt c{}_{n}^{(q)} \}_{n=0}^{+\infty}$.


Для этого сначала образуем промежуточную последовательность 


$\{d_{n}^{(q)} \}_{n=0}^{+\infty}$, положив


\begin{alignat*}{2}


d_{\nu}^{(q)}&=0 &\quad & \text{при }\ 0 \le \nu \le n_1^{(q)},\\


%


d_{\nu}^{(q)}&=(\min_{n_r^{(q)}<n<n_{r+1}^{(q)}} w_n) / n_{r+1}^{(q)} 


&\quad &\text{при } \ n_r^{(q)}< \nu \le n_{r+1}^{(q)}.


\end{alignat*}


Обозначим


$H_{ir}^{(q)}=\sum\limits_{\nu=l_{ir}^{(q)}+1}^{m_{ir}^{(q)}} d_{\nu}$ 


($i=1, \dotsc, k+2$)


и 


$H_{r}^{(q)}=\min\limits_{i=1,\dotsc, k+2} H_{ir}^{(q)}$.





Положим теперь


$\wt c{}_{\nu}^{(q)}=0$ при $0 \le \nu \le n_1^{(q)}$,


$\wt c{}_{\nu}^{(q)}=\frac{H_{r}^{(q)}}{H_{ir}^{(q)}} d_{\nu}^{(q)}$


при $l_{ir}^{(q)}<\nu \le m_{ir}^{(q)}$ ($i=1, \dotsc, k+2$),


$\wt c{}_{\nu}^{(q)}=0$ при $\nu \in 


(n_r^{(q)}; n_{r+1}^{(q)}] \setminus 


\mathop{\cup}\limits_{i=1}^{k+2} (l_{ir}^{(q)}; m_{ir}^{(q)}]$.


Для последовательности


$\{\wt c{}_{n}^{(q)} \}_{n=0}^{+\infty}$ выполнено


$0 \le \wt c{}_{n}^{(q)} \le d_n \le \frac{w_n}{n}=c_n$


и 


\begin{equation}


\sum_{\nu=l_{ir}^{(q)}+1}^{m_{ir}^{(q)}} \wt c{}_{\nu}^{(q)}= 


\frac{H_{r}^{(q)}}{H_{ir}^{(q)}} 


\sum_{\nu=l_{ir}^{(q)}+1}^{m_{ir}^{(q)}} d_{\nu}=H_r^{(q)}  


\ \ (i=1, \dotsc, k+2).


\end{equation}


Из равенства (35) вытекает, что выполнена формула (2) из 


определения условий $(ND, k)$. 





Поскольку $\lim\limits_{n \to+\infty} w_n=+\infty$,


то существует $N$ такое, что для всех $n>N$ выполнено


$$


w_n>\frac{2}{(q^p-1)^{k+1}}.


$$





Пусть $r$ любое такое, что $ n_r>N$. 


Тогда


\begin{multline*}


\displaybreak[3]


\sum_{i=1}^{k+2}\sum_{\nu=l_{ir}^{(q)}+1}^{m_{ir}^{(q)}} 


(n_{r+1}^{(q)}-\nu+1)^k \wt c{}_{\nu}^{(q)}=\sum_{i=1}^{k+2}


\frac{H_{r}^{(q)}}{H_{ir}^{(q)}} 


\sum_{\nu=l_{ir}^{(q)}+1}^{m_{ir}^{(q)}} 


(n_{r+1}^{(q)}-\nu+1)^k d_{\nu}=\\


%


=\frac{\min\limits_{n_r^{(q)}<n \le n_{r+1}^{(q)}} w_n}{n_{r+1}^{(q)}} 


\sum_{i=1}^{k+2}\frac{H_{r}^{(q)}}{H_{ir}^{(q)}} 


\sum_{\nu=l_{ir}^{(q)}+1}^{m_{ir}^{(q)}} 


(n_{r+1}^{(q)}-\nu+1)^k \ge \\


%


\ge \frac{2}{(q^p-1)^{k+1}}\frac{1}{n_{r+1}^{(q)}}\sum_{i=1}^{k+2}


\frac{H_{r}^{(q)}}{H_{ir}^{(q)}} 


\sum_{\nu=l_{ir}^{(q)}+1}^{m_{ir}^{(q)}} 


(n_{r+1}^{(q)}-m_{(k+2)r}^{(q)})^k=\\


%


= \frac{2}{(q^p-1)^{k+1}}


\frac{(m_{(k+2)r}^{(q)})^k}{n_{r+1}^{(q)}}


\bigg(\frac{n_{r+1}^{(q)}}{m_{(k+2)r}^{(q)}}-1 \bigg)^k


\ \sum_{i=1}^{k+2}\frac{H_{r}^{(q)}}{H_{ir}^{(q)}} 


l_{ir}^{(q)}\bigg(\frac{m_{ir}^{(q)}}{l_{ir}^{(q)}}-1 \bigg) \ge \\


%


\ge (n_r^{(q)})^k \sum_{i=1}^{k+2}


\frac{H_{r}^{(q)}}{H_{ir}^{(q)}} \ge


(n_r^{(q)})^k=\varepsilon (n_r^{(q)})^k.


\end{multline*}


Тем самым справедлива и формула (3) из определения условий $(ND, k)$. 


Следовательно,


$$


c_n \in (ND, k).\quad\square


$$





Автор глубоко признателен своему научному руководителю 


член-корреспонденту РАН\linebreak П.Л.\,Ульянову за внимание к работе и 


ряд полезных замечаний.





\begin{thebibliography}{9}





\bibitem{1}


Харди Г. {\em Расходящиеся ряды}. -- М.: Ин. лит., 


1951. -- 504~c.


\bibitem{2}


Hardy G.H. {\em Theorems relating to the summability 


and convergence of slowly oscillating series} // Proc. London Math. 


Soc. -- 1910. -- V.\,8. -- \No\,2. -- P.\,301--320.


\bibitem{3}


Littlewood J.E. {\em The converse of Abel's theorem 


on power series} // Proc. London Math. 


Soc. -- 1910. -- V.\,9. -- \No\,2. -- P.\,434--448.


\bibitem{4}


Lorentz G.G. {\em Tauberian theorems and tauberian conditions} // Trans. 


Amer. Math. Soc. -- 1948. -- V.\,63. -- \No\,2. -- P.\,226--234.


\bibitem{5}


Мельник В.И. {\em О суммировании рядов методами Чезаро и Абеля--Пуассона}


// Матем. сб. -- 1965. -- Т.\,67. -- \No\,4. -- С.\,535--540.


\bibitem{6}


Степанянц С.А. {\em Условия тауберова типа, связывающие 


методы суммирования Чезаро и Рисса различных порядков} // Anal.


Math. -- 2001. -- Т.\,27. -- С.\,287-318. 





\end{thebibliography}





\vskip 2cm





\post{\it Московский государственный}{\it Поступила}


\post{\it университет}{29.04.2003}





\label{end}


\end{document}


